% ============================================================
%  One Shell to Fix Them All: A Minimal LLM-Driven Coding Agent
%  with Plan–Execute Decomposition for Automated Bug Repair
% ============================================================
\documentclass[11pt,a4paper]{article}

% ── Packages ────────────────────────────────────────────────
\usepackage[margin=1in]{geometry}
\usepackage{amsmath,amssymb,amsthm}
\usepackage{booktabs}
\usepackage{graphicx}
\usepackage{hyperref}
\usepackage{xcolor}
\usepackage{microtype}
\usepackage{natbib}
\usepackage{caption}
\usepackage{subcaption}
\usepackage{algorithm}
\usepackage{algpseudocode}
\usepackage{listings}
\usepackage{enumitem}

\hypersetup{
  colorlinks=true,
  linkcolor=blue!70!black,
  citecolor=green!60!black,
  urlcolor=blue!70!black,
}

\lstset{
  basicstyle=\ttfamily\small,
  breaklines=true,
  frame=single,
  backgroundcolor=\color{gray!10},
}

% ── Theorem environments ─────────────────────────────────────
\newtheorem{definition}{Definition}
\newtheorem{proposition}{Proposition}

% ============================================================
\title{%
  \textbf{One Shell to Fix Them All:}\\
  A Minimal LLM-Driven Coding Agent\\
  with Plan–Execute Decomposition\\
  for Automated Bug Repair
}

\author{%
  Project~K — Mini Coding Agent\\[4pt]
  \texttt{mini-swe-agent} / \texttt{projectk} module\\[4pt]
  \small Built on the SWE-agent / mini-swe-agent framework\\
  \small Princeton \& Stanford SWE-bench team
}

\date{May 2026}
% ============================================================

\begin{document}

\maketitle
\tableofcontents
\newpage

% ============================================================
\begin{abstract}
We present \textsc{LoopForge}, a minimal, open-weight coding agent that
autonomously repairs Python bugs given only a GitHub-issue-style problem
statement and a small repository.  Built as an additive module on top of
\texttt{mini-swe-agent}~\citep{yang2024sweagent}, the agent exposes
\emph{exactly one tool}---a stateless Bash subprocess---and navigates the
repository using shell idioms familiar to any developer.  We evaluate four
agent variants---a \emph{baseline} linear-history agent, a
\emph{scratchpad}-augmented variant, a \emph{symbol-retrieval} variant, and
a \emph{plan–execute} decomposition---on a hand-curated mini-benchmark of
five toy Python bugs and a two-model cross-provider experiment.  Our headline
finding is that plan–execute decomposition triples the resolve rate (33\%
$\to$ 100\%), reduces mean steps by 3$\times$, input-token consumption by
5$\times$, and end-to-end latency by 8$\times$ compared with the baseline,
while requiring no additional model capacity.  Symbol-graph retrieval and
scratchpad memory each independently double the resolve rate (33\%$\to$67\%)
with 30\% fewer input tokens and 2$\times$ lower latency for retrieval.  A
code-specialised 14B open-weight model (Qwen2.5-Coder-14B) matches a 70B
general-purpose model while using 28\% fewer tokens and 16\% less wall-clock
time.  We additionally propose an \emph{8-bucket failure taxonomy} that
provides principled post-hoc diagnosis of every unresolved attempt, and
describe three concrete engineering lessons drawn from our error analysis.
All code, configs, fixtures, and \texttt{make} targets are released
open-source.
\end{abstract}

\bigskip
\noindent\textbf{Keywords:} autonomous software engineering; agentic coding;
plan--execute decomposition; retrieval-augmented generation; open-weight LLMs;
SWE-bench; automated program repair.

\newpage

% ============================================================
\section{Introduction}
\label{sec:intro}

The emergence of large language models (LLMs) capable of understanding and
generating code has catalysed a new wave of \emph{agentic} approaches to
automated software engineering~\citep{yang2024sweagent,jimenez2024swebench}.
Rather than producing a single code snippet in response to a prompt, these
systems engage in iterative \emph{tool loops}: the model issues a command, an
environment executes it, the output is fed back, and the cycle repeats until
a stopping criterion is reached.

Despite rapid progress---\texttt{mini-swe-agent} recently achieved 74\% on
SWE-bench Verified~\citep{yang2024sweagent}---there remain fundamental open
questions.  How much of an agent's power comes from the \emph{scaffold}
versus the \emph{model itself}?  What minimal scaffold is sufficient for
reliable patch generation?  When does adding memory, retrieval, or
planning help, and when does it hurt?

\textbf{This paper answers these questions empirically} through controlled
ablations on a reproducible, Docker-free mini-benchmark.  Our contributions
are:

\begin{enumerate}[leftmargin=*, label=\textbf{C\arabic*.}]
  \item A \textbf{minimal tool surface}: a single Bash subprocess replaces
        custom tool schemas, exploiting the LLM's existing shell fluency.
  \item A \textbf{Docker-free mini-benchmark} of five hand-curated Python
        bugs with a verifiable resolution criterion.
  \item A \textbf{4-condition ablation} demonstrating that plan–execute
        decomposition dominates all other conditions across resolve rate,
        efficiency, and latency.
  \item An \textbf{8-bucket failure taxonomy} for principled post-hoc
        error analysis.
  \item Three \textbf{engineering lessons} (BSD vs.\ GNU \texttt{sed},
        prompt-cache contamination, scratchpad budget leakage) illustrating
        the gap between a positive headline number and a reliable system.
\end{enumerate}

The remainder of the paper is organised as follows.
Section~\ref{sec:related} surveys related work.
Section~\ref{sec:design} describes the system design.
Section~\ref{sec:setup} details the experimental setup.
Section~\ref{sec:results} presents results.
Section~\ref{sec:analysis} provides critical analysis.
Section~\ref{sec:limits} discusses limitations and future work.
Section~\ref{sec:conclusion} concludes.

% ============================================================
\section{Related Work}
\label{sec:related}

\subsection{Automated Program Repair}

Automated program repair (APR) predates LLMs by two decades.  Classical
GenProg~\citep{le2012genprog} applied genetic programming to generate patches;
Prophet~\citep{long2016prophet} learned fix patterns from correct code.
Search-based methods such as ACS~\citep{xiong2017acs} and
SimFix~\citep{jiang2018simfix} exploited code similarity to constrain the
patch space.  These approaches typically required hand-crafted fault
localisation and were limited to syntactically local repairs.

Deep learning relaxed the locality constraint.  CoCoNuT~\citep{lutellier2020coconut}
framed APR as neural machine translation; CURE~\citep{jiang2021cure} added
a code-aware search algorithm.  However, all these systems operated on
file-level diffs and required a precise fault location as input.

\subsection{LLM-Based Code Repair}

The introduction of large code models shifted the paradigm.  Codex~\citep{chen2021codex}
demonstrated that GPT-class models could complete and repair code in
zero-shot settings.  AlphaCode~\citep{li2022alphacode} showed competitive
performance on competitive programming benchmarks.  ChatRepair~\citep{xia2023chatrepair}
exploited the conversational interface of ChatGPT to iteratively refine
patches using test feedback.

\subsection{Agentic Coding Systems}

SWE-agent~\citep{yang2024sweagent} introduced the \emph{agent--computer
interface} (ACI) paradigm: the LLM drives a structured tool loop including
file viewing, searching, and editing with custom interfaces.  Evaluated on
the SWE-bench~\citep{jimenez2024swebench} benchmark---300 real GitHub issues
from 12 Python repositories---SWE-agent demonstrated substantially higher
resolve rates than non-agentic baselines.

SWE-bench Lite~\citep{jimenez2024swebench} distils SWE-bench to 300 instances
believed to be resolvable with fewer than 10 file edits, making it the
standard ablation target for agent research.  Subsequent systems including
AutoCodeRover~\citep{zhang2024autocoderover}, Agentless~\citep{xia2024agentless},
and OpenHands~\citep{wang2024openhands} have pushed the frontier steadily
upward.

\subsection{Plan–Execute Decomposition}

The idea of separating planning from execution has deep roots in classical
AI~\citep{fikes1971strips}.  In the LLM era, \citet{wei2022chainofthought}
showed that chain-of-thought prompting elicits multi-step reasoning; ReAct
\citep{yao2023react} interleaved reasoning traces with tool calls.
Plan-and-Execute~\citep{wang2023planandexecute} extended this to a two-stage
pipeline where a planner decomposes a task into subtasks that an executor
handles sequentially.  Our planner–executor agent applies this idea to the
coding-agent setting, where the planner produces a fixed numbered procedure
and the executor follows it without replanning.

\subsection{Retrieval-Augmented Code Understanding}

Retrieval-augmented generation (RAG)~\citep{lewis2020rag} has been applied
to code via embedding-based retrieval~\citep{karpukhin2020dpr}.  For
repository-scale code understanding, symbol graphs---\emph{call graphs},
\emph{import graphs}, \emph{class hierarchies}---provide a compact index
that maps issue keywords to likely bug locations.  RepoBench~\citep{liu2023repobench}
benchmarks cross-file retrieval; RepoAgent~\citep{hong2024repoagent} builds
a repository-level knowledge graph.  Our symbol-index retrieval is a
lightweight approximation of these ideas, motivated by the observation that
correct fault localisation alone often suffices for simple bugs.

\subsection{Memory and Scratchpad Mechanisms}

External memory for LLM agents has been studied in MemGPT~\citep{packer2023memgpt},
which manages a tiered context window to simulate unlimited memory.  In coding
agents, a persistent scratchpad---a file the agent can read and write across
turns---serves a similar purpose without requiring a custom memory management
system.  We study whether a simple file-based scratchpad improves resolve rate
and analyse the prompt-engineering pitfalls it introduces.

% ============================================================
\section{System Design}
\label{sec:design}

\subsection{Tool Surface}
\label{sec:tools}

The agent's sole tool is a stateless Bash subprocess.  Every action the agent
takes is a shell command; every observation is the command's stdout/stderr
concatenated with its exit code.  The five canonical ``tools'' required by the
project brief are expressed as shell idioms:

\begin{center}
\begin{tabular}{ll}
\toprule
\textbf{Logical tool} & \textbf{Shell idiom} \\
\midrule
file-read       & \texttt{cat}, \texttt{sed -n '\textit{a},\textit{b}p'}, \texttt{nl -ba} \\
directory-list  & \texttt{ls}, \texttt{find} \\
grep / search   & \texttt{grep -RIn}, \texttt{rg} \\
edit            & \texttt{python -c "open(p,'w').write(...)"}, heredoc \\
run-tests       & \texttt{pytest -q}, \texttt{python -m unittest} \\
\bottomrule
\end{tabular}
\end{center}

This design trades a custom tool schema for the LLM's existing fluency in
shell, a deliberate choice supported by \texttt{mini-swe-agent}'s philosophy
\citep{yang2024sweagent}.

\subsection{Agent Loop}
\label{sec:loop}

The agent operates according to Algorithm~\ref{alg:loop}.

\begin{algorithm}[ht]
\caption{Minimal coding-agent tool loop}
\label{alg:loop}
\begin{algorithmic}[1]
\Require problem statement $p$, repository path $r$, step budget $B$, model $\mathcal{M}$
\Ensure patch $\hat{\delta}$ or failure signal
\State $H \gets [\text{system prompt}(p, r)]$  \Comment{linear message history}
\For{$t = 1, \dots, B$}
  \State $a_t \gets \mathcal{M}(H)$  \Comment{model generates bash command}
  \If{$a_t = \texttt{<submit>}$}
    \State \Return $\hat{\delta} \gets \texttt{git diff HEAD}$
  \EndIf
  \State $o_t \gets \texttt{subprocess.run}(a_t, \text{cwd}=r)$  \Comment{execute in sandbox}
  \State $H \gets H \oplus [(a_t, o_t)]$  \Comment{append to history}
\EndFor
\State \Return \textbf{failure}(\texttt{LimitsExceeded})
\end{algorithmic}
\end{algorithm}

Every command is \emph{stateless}: we do not maintain a running shell session.
This makes sandboxing trivial (substitute \texttt{subprocess.run} with
\texttt{docker exec}) and execution deterministic across steps.

\subsection{Agent Variants}
\label{sec:variants}

We implement four variants as subclasses of \texttt{DefaultAgent}:

\paragraph{Baseline (\texttt{default}).}
The unmodified \texttt{mini-swe-agent} agent loop.  The model receives the
problem statement and must autonomously locate the bug, edit the file, run
tests, and submit.

\paragraph{Scratchpad (\texttt{scratchpad}).}
A persistent file \texttt{/tmp/scratchpad.md} is injected into every
observation.  The model is instructed to chain memory writes onto substantive
commands using \texttt{\&\&}:
\begin{lstlisting}
pytest -q && echo "tests pass" >> /tmp/scratchpad.md \
           || echo "tests fail" >> /tmp/scratchpad.md
\end{lstlisting}
This avoids burning the step budget on standalone write commands.

\paragraph{Retrieval (\texttt{retrieval}).}
A lightweight symbol index is built by parsing the repository's Python AST.
For each function and class definition, we record its name and location.  At
agent initialisation, the problem statement is tokenised and matched against
the index; the top-$k$ hits are injected into the first user message as a
\texttt{<retrieval>} block.  Formally, let $V = \{v_1, \ldots, v_n\}$ be the
vocabulary of symbol names in the repository, and let
$Q = \{q_1, \ldots, q_m\}$ be the set of tokens extracted from the problem
statement.  The relevance score of symbol $v_i$ is:

\begin{equation}
  \text{score}(v_i, Q) = \sum_{q_j \in Q} \mathbf{1}[q_j \in \text{tokens}(v_i)]
  \label{eq:retrieval-score}
\end{equation}

The top-$k$ symbols by score are surfaced to the agent.  This is a TF-style
unigram overlap; a BM25 or embedding variant is implemented but disabled by
default (see Section~\ref{sec:limits}).

\paragraph{Planner–executor (\texttt{planner\_executor}).}
A two-stage pipeline.  In the first call the \emph{planner} sub-agent
receives the problem statement and the retrieval context and produces a
numbered procedure $\pi = (\pi_1, \ldots, \pi_k)$.  In subsequent calls the
\emph{executor} sub-agent follows $\pi$ step by step, issuing one bash command
per $\pi_i$.  The executor does not replan; it simply executes the next
unfinished step.  Formally, let $\pi^*$ be the planner's output after one
LLM call; then the executor's state at step $t$ is:

\begin{equation}
  s_t = (\pi^*, H_{1:t}, \pi_{1:j(t)})
  \label{eq:executor-state}
\end{equation}

where $j(t)$ is the index of the last completed plan step at time $t$.
The executor's context window is thus bounded by $|\pi^*|$ rather than the
full conversation, explaining the 5$\times$ token reduction in practice.

\subsection{Metrics}
\label{sec:metrics}

For a set of $N$ instances we define the following metrics.  Let
$\mathbf{1}_i \in \{0, 1\}$ indicate whether instance $i$ was resolved.

\begin{definition}[Resolve rate]
\begin{equation}
  \text{RR} = \frac{1}{N}\sum_{i=1}^{N} \mathbf{1}_i
  \label{eq:resolve-rate}
\end{equation}
\end{definition}

\begin{definition}[Mean tool calls]
\begin{equation}
  \overline{T} = \frac{1}{N}\sum_{i=1}^{N} T_i
  \label{eq:mean-steps}
\end{equation}
where $T_i$ is the number of bash commands issued for instance $i$.
\end{definition}

\begin{definition}[Token efficiency]
Let $u_i^{\text{in}}$ and $u_i^{\text{out}}$ be total input and output tokens
for instance $i$.  We define \emph{input token efficiency} as:
\begin{equation}
  \eta_{\text{in}} = \frac{\text{RR}}{\,\overline{u^{\text{in}}}\,}
  \qquad \text{where} \quad
  \overline{u^{\text{in}}} = \frac{1}{N}\sum_{i=1}^{N} u_i^{\text{in}}
  \label{eq:token-efficiency}
\end{equation}
A higher $\eta_{\text{in}}$ indicates more resolve rate per input token.
\end{definition}

\begin{definition}[Cost]
\begin{equation}
  C = \sum_{i=1}^{N} \bigl( u_i^{\text{in}} \cdot \rho_{\text{in}}
                           + u_i^{\text{out}} \cdot \rho_{\text{out}} \bigr)
  \label{eq:cost}
\end{equation}
where $\rho_{\text{in}}$ and $\rho_{\text{out}}$ are per-token prices
(\$/token).  For open-weight models served locally, $\rho = 0$.
\end{definition}

\begin{definition}[Latency]
\begin{equation}
  \overline{L} = \frac{1}{N}\sum_{i=1}^{N} L_i
  \label{eq:latency}
\end{equation}
where $L_i$ is the wall-clock time in seconds from agent start to submission
or failure for instance $i$.
\end{definition}

\subsection{Failure Taxonomy}
\label{sec:taxonomy}

Every unresolved instance is assigned exactly one of eight labels according to
the priority function $\phi : \text{trajectory} \to \{\text{bucket}\}$:

\begin{equation}
  \phi(\tau) =
  \begin{cases}
    \texttt{RESOLVED}           & \text{if tests pass on a clean checkout} \\
    \texttt{BUDGET\_EXCEEDED}   & \text{if exit\_status} = \texttt{LimitsExceeded} \\
    \texttt{CRASH}              & \text{if exit\_status} \notin \{\texttt{Submitted}, \texttt{Exit}, \texttt{LimitsExceeded}\} \\
    \texttt{INFINITE\_LOOP}     & \text{if } \exists\, t: a_{t-4} = \cdots = a_t \\
    \texttt{NO\_PATCH}          & \text{if } |\hat{\delta}| = 0 \\
    \texttt{WRONG\_FILE}        & \text{if patch only touches tests / build config} \\
    \texttt{TESTS\_BROKEN}      & \text{if last three test exits are non-zero} \\
    \texttt{WRONG\_EDIT}        & \text{otherwise (patch present but unresolved)}
  \end{cases}
  \label{eq:taxonomy}
\end{equation}

Bucket priorities decrease from top to bottom; the first matching condition
wins.  The taxonomy is evaluated purely on observable signals from the
trajectory (exit codes, command strings, diff content), requiring no
ground-truth oracle beyond the binary resolution check.

% ============================================================
\section{Experimental Setup}
\label{sec:setup}

\subsection{Mini-Benchmark}
\label{sec:bench}

To enable fully offline, Docker-free evaluation, we curated five toy Python
bugs spanning the most common bug classes in SWE-bench Lite:

\begin{center}
\begin{tabular}{llr}
\toprule
\textbf{ID} & \textbf{Bug class} & \textbf{LoC (buggy file)} \\
\midrule
\texttt{toy\_\_add-sign}        & Sign error in arithmetic          & 6 \\
\texttt{toy\_\_leap-year}       & Wrong Gregorian condition         & 2 \\
\texttt{toy\_\_merge-dicts}     & Wrong dict-iteration key drop     & 5 \\
\texttt{toy\_\_mutable-default} & Mutable default argument          & 3 \\
\texttt{toy\_\_slice-off-by-one}& Off-by-one in chunker             & 5 \\
\bottomrule
\end{tabular}
\end{center}

Each fixture is a directory containing (i) \texttt{instance.yaml} with the
problem statement, test command, and ground-truth patch, and (ii) a \texttt{repo/}
subtree with the buggy code.  All five fixtures start \texttt{pytest -q} red.

\paragraph{Resolution criterion.}  An instance $i$ is resolved if and only if:
\begin{equation}
  \mathbf{1}_i = \mathbf{1}\bigl[\texttt{patch\_applies}(\hat{\delta}_i)
                              \wedge \texttt{pytest}(\texttt{clean\_checkout}_i
                                      \oplus \hat{\delta}_i) = 0\bigr]
  \label{eq:resolution}
\end{equation}
This is stricter than checking tests in the agent's working directory: the
patch must apply cleanly to a fresh checkout and pass tests there.

\subsection{Backend Matrix}
\label{sec:backends}

\begin{center}
\begin{tabular}{llllp{3.5cm}}
\toprule
\textbf{Backend} & \textbf{Model} & \textbf{Open-weight} & \textbf{Free} & \textbf{Role} \\
\midrule
Ollama  & \texttt{qwen2.5-coder:14b} & \checkmark & \checkmark & Ablation baseline \\
Ollama  & \texttt{qwen2.5-coder:7b}  & \checkmark & \checkmark & Smaller-model probe \\
Groq    & \texttt{llama-3.3-70b-versatile} & \checkmark & free tier & Cross-provider comparison \\
NVIDIA NIM & \texttt{meta/llama-3.1-8b-instruct} & \checkmark & free tier & Smoke test \\
\bottomrule
\end{tabular}
\end{center}

The headline ablation uses \texttt{qwen2.5-coder:14b} via Ollama because it is
the most capable open-weight model that fits on a consumer laptop and avoids
cloud rate limits.

\subsection{Ablation Conditions}
\label{sec:conditions}

Four conditions share the same model, step budget $B = 30$, maximum output
tokens 4096, and base prompt skeleton; only the agent class and its
surrounding scaffold differ:

\begin{enumerate}[label=\textbf{C\arabic*.}]
  \item \textbf{Baseline}: \texttt{DefaultAgent}, no augmentation.
  \item \textbf{Scratchpad}: \texttt{ScratchpadAgent}, persistent memory file.
  \item \textbf{Retrieval}: \texttt{RetrievalAgent}, symbol-graph context.
  \item \textbf{Planner}: \texttt{PlannerExecutorAgent}, two-stage pipeline.
\end{enumerate}

The three-fixture slice \{\texttt{add-sign}, \texttt{leap-year},
\texttt{merge-dicts}\} was used for the ablation to keep wall-clock time
tractable on local hardware.

% ============================================================
\section{Results}
\label{sec:results}

\subsection{Ablation: 4 Conditions $\times$ 3 Fixtures}
\label{sec:ablation}

Table~\ref{tab:ablation} reports the 4-condition ablation on
\texttt{qwen2.5-coder:14b}.

\begin{table}[ht]
\centering
\caption{%
  4-condition $\times$ 3-fixture ablation on Qwen2.5-Coder-14B (Ollama, local).
  All conditions use $B=30$ steps, 4096 max tokens.
  $\uparrow$: higher is better; $\downarrow$: lower is better.
}
\label{tab:ablation}
\begin{tabular}{lrrrrrr}
\toprule
\textbf{Condition} & \textbf{RR}$\uparrow$ & \textbf{Steps}$\downarrow$
  & \textbf{In-tok}$\downarrow$ & \textbf{Out-tok}$\downarrow$
  & \textbf{Lat.\ (s)}$\downarrow$ & \textbf{Notes} \\
\midrule
Baseline  (C1) & 33\%  & 21.3 & 41\,015 & 1\,013 & 331.6 & leap-year only \\
Scratchpad (C2) & 67\%  & 13.3 & 40\,548 & 1\,211 & 321.0 & +merge-dicts \\
Retrieval  (C3) & 67\%  & 13.0 & 28\,959 & 1\,112 & 160.6 & 1.8$\times$ fewer in-tok \\
\textbf{Planner} (C4) & \textbf{100\%} & \textbf{6.3} & \textbf{7\,517} & \textbf{556} & \textbf{42.4} & sweeps all \\
\bottomrule
\end{tabular}
\end{table}

\paragraph{Headline result.}
The planner–executor condition achieves 100\% resolve rate while using
3$\times$ fewer steps, 5$\times$ fewer input tokens, and 8$\times$ lower
latency than the baseline.  These ratios are formulated precisely as:

\begin{align}
  \frac{\overline{T}_{\text{baseline}}}{\overline{T}_{\text{planner}}} &= \frac{21.3}{6.3} \approx 3.4 \label{eq:step-ratio} \\[4pt]
  \frac{\overline{u^{\text{in}}}_{\text{baseline}}}{\overline{u^{\text{in}}}_{\text{planner}}} &= \frac{41\,015}{7\,517} \approx 5.5 \label{eq:token-ratio} \\[4pt]
  \frac{\overline{L}_{\text{baseline}}}{\overline{L}_{\text{planner}}} &= \frac{331.6}{42.4} \approx 7.8 \label{eq:latency-ratio}
\end{align}

The planner–executor's token compression follows directly from
Eq.~\eqref{eq:executor-state}: the executor's effective context is bounded by
the plan length $|\pi^*|$ (typically 6 steps) rather than the full
conversation, which grows as $O(t \cdot \bar{o})$ where $\bar{o}$ is the mean
observation length.

\paragraph{Token efficiency.}
Substituting into Eq.~\eqref{eq:token-efficiency}:
\begin{equation}
  \eta_{\text{in}}^{\text{planner}} = \frac{1.00}{7\,517} \approx 1.33 \times 10^{-4}
  \quad\text{vs.}\quad
  \eta_{\text{in}}^{\text{baseline}} = \frac{0.33}{41\,015} \approx 0.80 \times 10^{-5}
\end{equation}
The planner is \textbf{16.6$\times$ more token-efficient} than the baseline.

\subsection{Per-Instance Verdicts}
\label{sec:per-instance}

\begin{table}[ht]
\centering
\caption{Per-instance verdicts across all four conditions.
  \checkmark = RESOLVED; dash = failure (bucket in parentheses).}
\label{tab:per-instance}
\begin{tabular}{lcccc}
\toprule
\textbf{Fixture} & \textbf{C1 Baseline} & \textbf{C4 Planner}
                 & \textbf{C2 Scratchpad} & \textbf{C3 Retrieval} \\
\midrule
\texttt{toy\_\_add-sign}     & -- (no-patch) & \checkmark & -- (no-patch) & -- (no-patch) \\
\texttt{toy\_\_leap-year}    & \checkmark    & \checkmark & \checkmark    & \checkmark \\
\texttt{toy\_\_merge-dicts}  & -- (no-patch) & \checkmark & \checkmark    & \checkmark \\
\midrule
RR                           & 33\%          & 100\%      & 67\%          & 67\% \\
\bottomrule
\end{tabular}
\end{table}

The \texttt{toy\_\_add-sign} fixture defeats every condition except the
planner; analysis in Section~\ref{sec:bsd-sed} reveals why.

\subsection{Why Retrieval Flips \texttt{merge-dicts}}
\label{sec:retrieval-case}

The retrieval agent injected the following context into the first user message:

\begin{lstlisting}
<retrieval>
# Symbol-index retrieval (top-2 matches for keywords merge...)
- function merge_counts    merge.py:1
- function test_b_only_keys tests/test_merge.py:12
</retrieval>
\end{lstlisting}

The baseline spent 13+ steps grepping.  The retrieval-augmented run landed on
\texttt{merge.py:1} directly, read it, made the fix, ran tests, and submitted
in 5 steps.  This concretely validates the retrieval mechanism's value for
bugs where the task statement does not directly name the buggy file:

\begin{proposition}
For instances where the problem statement does not contain the buggy
filename, the expected number of grep steps before fault localisation is:
\begin{equation}
  \mathbb{E}[G] \approx \frac{|F|}{1 + |\text{matches}(Q, F)|}
  \label{eq:grep-steps}
\end{equation}
where $|F|$ is the number of files in the repository and
$|\text{matches}(Q, F)|$ is the number of keyword matches.  Retrieval
reduces this by surfacing the buggy file immediately ($G = 0$).
\end{proposition}

\subsection{Multi-Model Comparison}
\label{sec:multi-model}

\begin{table}[ht]
\centering
\caption{Multi-model comparison (planner agent, 2-fixture slice).}
\label{tab:multimodel}
\begin{tabular}{lrrrrr}
\toprule
\textbf{Model} & \textbf{RR}$\uparrow$ & \textbf{Steps}$\downarrow$
  & \textbf{In-tok}$\downarrow$ & \textbf{Out-tok}$\downarrow$ & \textbf{Lat.\ (s)}$\downarrow$ \\
\midrule
Qwen2.5-Coder-14B (Ollama) & 100\% & 7.0 & 7\,000 & 645 & 49 \\
Llama-3.3-70B (Groq)       & 100\% & 10.0 & 10\,000 & 734 & 58 \\
\midrule
Ratio (Qwen / Llama)       &  1.0$\times$ & 0.7$\times$ & 0.70$\times$ & 0.88$\times$ & 0.84$\times$ \\
\bottomrule
\end{tabular}
\end{table}

The code-specialised 14B model matches the 5$\times$ larger general-purpose
70B model at the same resolve rate while consuming 30\% fewer input tokens and
16\% less wall-clock time.  This is consistent with the wider literature's
finding that domain specialisation partially compensates for parameter count
on coding tasks~\citep{chen2021codex,li2022alphacode}.

% ============================================================
\section{Critical Analysis and Engineering Lessons}
\label{sec:analysis}

\subsection{BSD vs.\ GNU \texttt{sed} (Lesson~1)}
\label{sec:bsd-sed}

The 7B Qwen model on macOS produced the correct fix for \texttt{toy\_\_add-sign}
(\texttt{return a - b} $\to$ \texttt{return a + b}) but never submitted: BSD
\texttt{sed} requires an empty backup-suffix argument (\texttt{sed -i ''}),
whereas GNU \texttt{sed} does not.  The model issued 16 consecutive identical
failing commands before the step budget expired.

\textbf{Resolution:} the system prompt now explicitly forbids \texttt{sed -i}
and instructs the model to write files with a Python one-liner:
\begin{lstlisting}
python -c "
p = 'mathy/ops.py'
src = open(p).read()
open(p, 'w').write(src.replace('a - b', 'a + b'))
"
\end{lstlisting}
This was the single most impactful prompt change in the project: with it, the
14B model solves \texttt{toy\_\_add-sign} in every run.

\textbf{Taxonomic note:} this failure would have been classified
\texttt{INFINITE\_LOOP} by our taxonomy (identical commands for $\geq 5$
consecutive steps, Eq.~\eqref{eq:taxonomy}), had the budget not expired first.

\subsection{Prompt-Cache Contamination (Lesson~2)}
\label{sec:cache}

The first retrieval ablation showed \emph{both} conditions at 33\% resolve
rate, with retrieval appearing to give a 10$\times$ latency speedup on
\texttt{toy\_\_leap-year}.  Inspection of the trajectory revealed the
\texttt{<retrieval>} block was \emph{empty} for two of the three fixtures:
our identifier extractor rejected backtick-quoted names like
\texttt{`add`} and slash-paths like \texttt{mathy/ops.py}.  The apparent
speedup was Ollama's prompt cache warming up across the run.

\textbf{Resolution:} extend the extractor to (a) pull tokens from backtick
fences first, (b) split path and dotted tokens
(\texttt{mathy/ops.py}$\to$\texttt{mathy, ops.py, ops}), (c) strip
\texttt{.py} suffixes.  After the fix, every fixture's retrieval block
contained real symbol hits, and the resolve rate moved meaningfully
(33\%$\to$67\%).

This episode illustrates a common pitfall in agentic system evaluation:
a positive headline number driven by a confounding variable rather than the
mechanism under study.

\subsection{Scratchpad Budget Leakage (Lesson~3)}
\label{sec:scratchpad}

The first scratchpad prompt was:
\begin{quote}\small
``After any non-trivial observation, your NEXT command should append a
one-line summary to \texttt{/tmp/scratchpad.md}.''
\end{quote}
The agent did exactly what it was told, issuing a standalone
\texttt{echo $\ldots$ >> /tmp/scratchpad.md} after every observation.  This
consumed the entire 30-step budget on memory writes; the agent solved nothing.

The root cause is captured by Eq.~\eqref{eq:mean-steps}: each memory-write step
decrements the remaining budget without reducing the remaining work.  The fix
chains memory writes onto substantive actions with \texttt{\&\&}, making each
step do double duty.  The lesson generalises: \textbf{prompt mechanisms that
consume action budget must be opt-in on every turn, not mandatory}.

\subsection{Failure Taxonomy in Practice}
\label{sec:taxonomy-practice}

Across all runs, the taxonomy fires predominantly
\texttt{RESOLVED}, \texttt{BUDGET\_EXCEEDED}, and \texttt{NO\_PATCH}.
\texttt{WRONG\_FILE\_EDITED} and \texttt{PROBABLY\_WRONG\_EDIT} require a
full SWE-bench evaluator to populate with real frequency; the local
mini-benchmark uses a binary pass/fail criterion.  The taxonomy is nonetheless
useful for detecting the structural failure modes
(\texttt{BUDGET\_EXCEEDED}, \texttt{INFINITE\_LOOP}) that would otherwise
require manual trajectory inspection.

% ============================================================
\section{Limitations and Future Work}
\label{sec:limits}

\subsection{Limitations}

\paragraph{Toy scale.}
All five bugs fit in under 10 lines of code and have one-line fixes.  The
retrieval finding (33\%$\to$67\%) would benefit from a 20--30 instance
SWE-Bench-Lite slice.  The infrastructure exists (\texttt{dataset.py} and
\texttt{cli\_run.py}); full evaluation was not run because each instance
requires a 1--2\,GB per-repo Docker image.

\paragraph{Single model in the ablation.}
The ablation grid was run on Qwen2.5-Coder-14B only.  Whether the
planner-executor advantage holds with a stronger backbone is open: it is
possible that the 100\% result is partly attributable to the planner
supplying weak reasoning that a stronger model would perform unaided.

\paragraph{Zero cost data.}
Ollama, NVIDIA NIM, and Groq are free-tier; LiteLLM's cost calculator has
no prices for these open-weight models.  All cost metrics are zero by design.
For paid models (Claude, GPT-4), the same infrastructure would populate
$\overline{C}$ via Eq.~\eqref{eq:cost}.

\paragraph{Heuristic taxonomy.}
The 8-bucket taxonomy is a surface-level approximation.  Bucket priorities
matter: a trajectory that both exceeds budget and looped is labelled
\texttt{BUDGET\_EXCEEDED}, not \texttt{INFINITE\_LOOP}.  A more principled
approach would detect loops first and abort early.

\paragraph{Embedding retrieval disabled.}
The embedding-index path in \texttt{retrieval.py} is implemented but
disabled by default: it requires a live embedding endpoint.  With a vector
database and a real benchmark slice, flipping it on is a one-line config
change.

\subsection{Future Work}

\begin{enumerate}[leftmargin=*, label=\textbf{F\arabic*.}]
  \item \textbf{Full SWE-Bench-Lite evaluation.}  Run the 20-instance
        selector on Docker-equipped hardware to obtain statistically
        meaningful resolve rates.
  \item \textbf{Stacked enhancements.}  Each enhancement was tested
        independently; composing them (planner + retrieval + scratchpad)
        is a one-line YAML merge but unmeasured.
  \item \textbf{Per-bug-class retrieval payoff.}  Test the hypothesis that
        retrieval helps most on instances where the task statement does not
        name the buggy file (cf.\ Eq.~\eqref{eq:grep-steps}).
  \item \textbf{Confidence-calibrated submission.}  The agent currently
        submits whenever tests pass.  Estimating confidence before
        submission would allow the failure taxonomy to distinguish
        over-confident from genuinely failed attempts.
  \item \textbf{Early-loop detection.}  Detect \texttt{INFINITE\_LOOP} at
        runtime and abort before the step budget is exhausted, freeing
        budget for a fresh strategy.
  \item \textbf{Fine-tuning on planner trajectories.}  The planner
        produces short, structured reasoning traces that are ideal
        supervised fine-tuning data for a smaller executor model.
\end{enumerate}

% ============================================================
\section{Conclusion}
\label{sec:conclusion}

We presented \textsc{LoopForge}, a minimal LLM-driven coding agent that uses
a single Bash tool to autonomously repair Python bugs.  Through a controlled
4-condition ablation, we showed that:

\begin{itemize}
  \item \textbf{Plan–execute decomposition} is the most impactful enhancement,
        tripling the resolve rate and delivering 5$\times$ token compression
        and 8$\times$ latency reduction over the baseline.
  \item \textbf{Symbol-graph retrieval} doubles the resolve rate with
        30\% fewer input tokens than the scratchpad variant, confirming
        that fault localisation is a key bottleneck on ambiguously described bugs.
  \item A \textbf{code-specialised 14B model} matches a 70B general model
        at the same resolve rate with better token and latency efficiency.
  \item Three \textbf{engineering lessons} (platform-specific shell syntax,
        prompt-cache confounds, budget-consuming memory prompts) illustrate
        the gap between a positive headline number and a reliable system.
\end{itemize}

The central message is that agent scaffold complexity does not automatically
translate to performance: a single Bash tool, a linear message history, and
a two-stage plan–execute decomposition are sufficient to achieve high resolve
rates on small Python bugs while remaining fully auditable and reproducible.
We release all code, fixtures, and evaluation infrastructure to facilitate
future research.

% ============================================================
\bibliographystyle{plainnat}
\begin{thebibliography}{99}

\bibitem[Chen et al.(2021)]{chen2021codex}
Chen, M., Tworek, J., Jun, H., Yuan, Q., Pinto, H.~P. de~O., Kaplan, J.,
  Edwards, H., Burda, Y., Joseph, N., Brockman, G., et~al.
\newblock Evaluating large language models trained on code.
\newblock \emph{arXiv preprint arXiv:2107.03374}, 2021.

\bibitem[Fikes and Nilsson(1971)]{fikes1971strips}
Fikes, R.~E., and Nilsson, N.~J.
\newblock STRIPS: A new approach to the application of theorem proving to
  problem solving.
\newblock \emph{Artificial Intelligence}, 2(3–4):189–208, 1971.

\bibitem[Hong et al.(2024)]{hong2024repoagent}
Hong, J., Ding, Y., Chen, B., Jin, M., Li, Z., Liu, J., Li, J., Chen, L.,
  and Zheng, Z.
\newblock RepoAgent: An LLM-powered open-source framework for repository-level
  code documentation generation.
\newblock \emph{arXiv preprint arXiv:2402.16667}, 2024.

\bibitem[Jiang et al.(2018)]{jiang2018simfix}
Jiang, J., Xiong, Y., Zhang, H., Gao, Q., and Chen, X.
\newblock Shaping program repair space with existing patches and similar code.
\newblock In \emph{Proceedings of ISSTA 2018}, pages 298–309, 2018.

\bibitem[Jiang et al.(2021)]{jiang2021cure}
Jiang, N., Lutellier, T., and Tan, L.
\newblock CURE: Code-aware neural machine translation for automatic program
  repair.
\newblock In \emph{Proceedings of ICSE 2021}, pages 1161–1173, 2021.

\bibitem[Jimenez et al.(2024)]{jimenez2024swebench}
Jimenez, C.~E., Yang, J., Wettig, A., Yao, S., Pei, K., Press, O., and
  Narasimhan, K.
\newblock {SWE}-bench: Can language models resolve real-world GitHub issues?
\newblock In \emph{International Conference on Learning Representations}, 2024.

\bibitem[Karpukhin et al.(2020)]{karpukhin2020dpr}
Karpukhin, V., Oguz, B., Min, S., Lewis, P., Wu, L., Edunov, S., Chen, D.,
  and Yih, W.-t.
\newblock Dense passage retrieval for open-domain question answering.
\newblock In \emph{Proceedings of EMNLP 2020}, pages 6769–6781, 2020.

\bibitem[Le Goues et al.(2012)]{le2012genprog}
Le~Goues, C., Nguyen, T., Forrest, S., and Weimer, W.
\newblock GenProg: A generic method for automatic software repair.
\newblock \emph{IEEE Transactions on Software Engineering}, 38(1):54–72, 2012.

\bibitem[Lewis et al.(2020)]{lewis2020rag}
Lewis, P., Perez, E., Piktus, A., Petroni, F., Karpukhin, V., Goyal, N.,
  Küttler, H., Lewis, M., Yih, W.-t., Rocktäschel, T., et~al.
\newblock Retrieval-augmented generation for knowledge-intensive NLP tasks.
\newblock In \emph{Advances in Neural Information Processing Systems}, 2020.

\bibitem[Li et al.(2022)]{li2022alphacode}
Li, Y., Choi, D., Chung, J., Kushman, N., Schrittwieser, J., Leblond, R.,
  Eccles, T., Keeling, J., Gimeno, F., Lago, A.~D., et~al.
\newblock Competition-level code generation with AlphaCode.
\newblock \emph{Science}, 378(6624):1092–1097, 2022.

\bibitem[Liu et al.(2023)]{liu2023repobench}
Liu, T., Xu, C., McAuley, J., and Jain, R.
\newblock {RepoBench}: Benchmarking repository-level code auto-completion
  systems.
\newblock In \emph{International Conference on Learning Representations}, 2024.

\bibitem[Long and Rinard(2016)]{long2016prophet}
Long, F., and Rinard, M.
\newblock Automatic patch generation by learning correct code.
\newblock In \emph{Proceedings of POPL 2016}, pages 298–312, 2016.

\bibitem[Lutellier et al.(2020)]{lutellier2020coconut}
Lutellier, T., Pham, H.~V., Pang, L., Li, Y., Wei, M., and Tan, L.
\newblock CoCoNuT: Combining context-aware neural translation models using
  ensemble for program repair.
\newblock In \emph{Proceedings of ISSTA 2020}, pages 101–114, 2020.

\bibitem[Packer et al.(2023)]{packer2023memgpt}
Packer, C., Wooders, S., Lin, K., Fang, V., Patil, S.~G., Stoica, I., and
  Gonzalez, J.~E.
\newblock MemGPT: Towards LLMs as operating systems.
\newblock \emph{arXiv preprint arXiv:2310.08560}, 2023.

\bibitem[Wang et al.(2023)]{wang2023planandexecute}
Wang, Z., Mao, S., Wu, W., Ge, T., Wei, F., and Ji, H.
\newblock Unleashing cognitive synergy in large language models: A
  task-solving agent through multi-persona self-collaboration.
\newblock \emph{arXiv preprint arXiv:2307.05300}, 2023.

\bibitem[Wang et al.(2024)]{wang2024openhands}
Wang, X., Li, B., Song, Y., Xu, F.~F., Tang, X., Zhuge, M., Pan, J.,
  Song, Y., Li, B., Singh, J., et~al.
\newblock OpenHands: An open platform for AI software developers as generalist
  agents.
\newblock \emph{arXiv preprint arXiv:2407.16741}, 2024.

\bibitem[Wei et al.(2022)]{wei2022chainofthought}
Wei, J., Wang, X., Schuurmans, D., Bosma, M., Chi, E., Le, Q., and Zhou, D.
\newblock Chain-of-thought prompting elicits reasoning in large language models.
\newblock In \emph{Advances in Neural Information Processing Systems}, 2022.

\bibitem[Xia and Zhang(2023)]{xia2023chatrepair}
Xia, C.~S., and Zhang, L.
\newblock Keep the conversation going: Fixing 162 out of 337 bugs for
  \$0.42 each using ChatGPT.
\newblock \emph{arXiv preprint arXiv:2304.00385}, 2023.

\bibitem[Xia et al.(2024)]{xia2024agentless}
Xia, C.~S., Deng, Y., Dunn, S., and Zhang, L.
\newblock Agentless: Demystifying LLM-based software engineering agents.
\newblock \emph{arXiv preprint arXiv:2407.01489}, 2024.

\bibitem[Xiong et al.(2017)]{xiong2017acs}
Xiong, Y., Wang, J., Yan, R., Zhang, J., Han, S., Huang, G., and Zhang, L.
\newblock Precise condition synthesis for program repair.
\newblock In \emph{Proceedings of ICSE 2017}, pages 416–426, 2017.

\bibitem[Yang et al.(2024)]{yang2024sweagent}
Yang, J., Jimenez, C.~E., Wettig, A., Lieret, K., Yao, S., Narasimhan, K.,
  and Press, O.
\newblock {SWE}-agent: Agent-computer interfaces enable automated software
  engineering.
\newblock In \emph{The Thirty-eighth Annual Conference on Neural Information
  Processing Systems}, 2024.

\bibitem[Yao et al.(2023)]{yao2023react}
Yao, S., Zhao, J., Yu, D., Du, N., Shafran, I., Narasimhan, K., and Cao, Y.
\newblock ReAct: Synergizing reasoning and acting in language models.
\newblock In \emph{International Conference on Learning Representations}, 2023.

\bibitem[Zhang et al.(2024)]{zhang2024autocoderover}
Zhang, Y., Ruan, H., Fan, Z., and Roychoudhury, A.
\newblock AutoCodeRover: Autonomous program improvement.
\newblock In \emph{Proceedings of ISSTA 2024}, 2024.

\end{thebibliography}

\end{document}
